\clearpage
\section{Problem 4: Atap Seng di Siang Bolong}
\label{sec:problem_4}

\subsection{Identifikasi Masalah}
Pemantauan kebakaran hutan dan lahan melalui instrumen satelit penginderaan jauh menghadirkan karakteristik beban yang sangat berbeda dibandingkan pemantauan seismik atau vulkanik. Satelit pengamat termal melintasi wilayah dalam orbit berkala dan memancarkan semburan data berukuran sangat besar (\textit{bursty data}) yang mencakup ribuan deteksi titik panas (\textit{hotspot}) dalam waktu singkat, kemudian hening hingga lintasan orbit berikutnya. Sebagian besar titik panas yang terdeteksi merupakan pantulan termal non-bahaya (seperti atap seng atau permukaan batuan terbuka yang terpanggang sinar matahari), sementara kebakaran hutan nyata sering kali teramati oleh beberapa instrumen satelit sekaligus pada waktu dan koordinat yang sedikit berbeda.

Pola integrasi lama berbasis tarikan berkala (\textit{polling}) tidak lagi memadai untuk menangani volume dan latensi aliran data ini. Beberapa persoalan mendasar yang dihadapi mencakup:
\begin{enumerate}
    \item \textbf{Pengelompokan Spasio-Temporal (\textit{Spatio-Temporal Event Aggregation}):} Satu titik panas individual belum merepresentasikan satu kejadian kebakaran hutan. Sistem harus mampu mendeteksi korelasi spasial (kedekatan koordinat geografis) dan korelasi temporal antar-deteksi dari berbagai satelit untuk diagregasikan menjadi satu entitas kejadian bahaya tunggal (\textit{HazardEvent WILDFIRE}) dengan atribut gabungan yang dinamis (seperti daya pancar radiasi kebakaran/FRP kumulatif dan tingkat keyakinan tertinggi).
    \item \textbf{Data Terlambat dan Ketidaksepakatan Jam Satelit (\textit{Out-of-Order and Delayed Events}):} Jam pada masing-masing instrumen satelit tidak selalu tersinkronisasi sempurna, dan transmisi data dari lintasan satelit dapat mengalami keterlambatan transfer jaringan hingga beberapa jam. Pengelompokan kejadian tidak boleh didasarkan pada waktu kedatangan data di server BNPB (\textit{processing time}), melainkan harus berbasis pada waktu observasi nyata satelit (\textit{event time}). Hal ini memunculkan dilema penutupan jendela agregasi: kapan sebuah jendela waktu boleh ditutup dan dianggap selesai jika data terkait masih berpotensi datang terlambat?
    \item \textbf{Ketahanan Aliran Data Berkecepatan Tinggi (\textit{Backpressure \& Burst Ingestion}):} Ketika aliran data dialirkan dengan akselerasi tinggi (misalnya simulasi percepatan 60$\times$), sistem harus mampu menyerap lonjakan tanpa kehilangan data (\textit{zero data loss}), membatasi degradasi latensi ($p95$), serta mengembalikan kondisi pemrosesan ke status mutakhir (\textit{near-zero lag}) secara cepat setelah lonjakan mereda.
    \item \textbf{Konsistensi Status Aliran Terhadap Kegagalan (\textit{Fault-Tolerant Stateful Stream Processing}):} Jika komponen pemroses mengalami pemadaman mendadak di tengah proses pengenalan kebakaran hutan, status sementara (\textit{intermediate state}) tidak boleh hilang atau rusak. Sistem yang pulih harus mampu merekonstruksi status agregasi secara identik persis dengan pemrosesan tanpa gangguan (\textit{exactly-once processing semantics}), menjamin tidak ada kebakaran yang terlewat atau terduplikasi.
\end{enumerate}

Dalam literatur sistem terdistribusi dan pemrosesan data beraliran (\textit{stream processing}), persoalan ini diidentifikasi sebagai \textbf{\textit{Event-Time vs. Processing-Time Windowing}}, \textbf{\textit{Watermarking and Allowed Lateness Handling}}, \textbf{\textit{Spatial Clustering over Streaming Data}}, \textbf{\textit{Stateful Stream Checkpointing}}, dan \textbf{\textit{Backpressure Management}}.

\subsection{Rancangan Solusi dan Mekanisme}
Untuk menangani semburan data titik panas satelit VIIRS secara presisi, tahan keterlambatan, dan konsisten terhadap kegagalan, Aggregator meninggalkan pola \textit{polling} M1 dan mengimplementasikan pipa pemrosesan aliran berkeadaan (\textit{stateful stream processing}) \cite{kleppmann2017designing,akidau2015dataflow}:

\begin{enumerate}
    \item \textbf{Klien Aliran SSE Berbasis Offset Monoton:} Klien SSE (\path{services/aggregator/internal/satellite/stream.go}) membuka koneksi persisten \texttt{GET} \path{/v1/stream?from=<offset>}, mengabaikan \textit{offset} yang sudah diterapkan, dan menyambung ulang apabila mendeteksi celah \textit{offset} alih-alih melompatinya. Respons HTTP~401, 403, dan 400 menghentikan aliran agar tidak memicu kesalahan berulang, sedangkan HTTP~429 dan 503 memicu \textit{backoff} yang menghormati \textit{header} \texttt{Retry-After}. Pesan \textit{heartbeat} yang membawa \texttt{id} turut memajukan posisi \textit{offset} karena memakai urutan \textit{offset} yang sama.
    \item \textbf{Identitas Deteksi Deterministik:} Kunci deteksi dibentuk dari gabungan teks mentah \path{satellite|acq_date|acq_time|latitude|longitude} persis seperti yang dikirim feed tanpa pembulatan numerik. Koordinat dalam notasi eksponen (misalnya \texttt{1.9e-4}) dinormalisasi menjadi desimal biasa (\texttt{0.00019}) melalui pergeseran titik desimal agar cocok dengan format kunci pada event \texttt{update} dan \texttt{retract}. Deteksi identik yang diterima berulang hanya dihitung satu kali.
    \item \textbf{Mesin Pengelompokan Inkremental (\textit{Connected Components}):} Mesin kebakaran (\path{services/aggregator/internal/satellite/engine.go}) mengindeks deteksi berdasarkan pita lintang selebar $0,01^\circ$ sehingga pencarian tetangga dalam radius $\le 1,0$~km dan selisih waktu kejadian $\le 360$~menit cukup memeriksa pita berdekatan. Himpunan komponen terhubung dengan $\ge 2$ deteksi membentuk satu kebakaran (\texttt{HazardEvent} \texttt{SATELIT}/\texttt{WILDFIRE}). Deteksi baru yang menjembatani beberapa komponen menggabungkan kebakaran terkait: kebakaran terbesar mempertahankan ID-nya sementara lainnya dipensiunkan (\texttt{retired\_reason = "merged"}). Pada peristiwa \texttt{retract} yang memecah komponen, potongan terbesar mempertahankan ID dan komponen yang menyusut menjadi $<2$ deteksi dipensiunkan (\texttt{below\_threshold}). ID kebakaran baru diturunkan dari SHA-256 kunci jangkar agar pemrosesan ulang menghasilkan ID identik. Deteksi ber-\textit{confidence} rendah (\texttt{l}) tetap disimpan tanpa membentuk kebakaran sehingga event \texttt{update} yang menaikkan \textit{confidence} tetap dapat diterapkan.
    \item \textbf{\textit{Event-Time Watermark} dan Penanganan Data Terlambat:} \textit{Watermark} dihitung sebagai waktu kejadian deteksi terbaru dikurangi \path{SATELLITE_LATENESS_MINUTES} (bawaan 360~menit). Deteksi terlambat yang masih berada dalam batas \textit{watermark} langsung memperbarui kebakaran terkait, sedangkan deteksi di luar batas disimpan terpisah dan dicatat pada metrik \texttt{late\_beyond\_watermark} tanpa dibuang diam-diam.
    \item \textbf{Pemetaan Kontrak dan Siklus Hidup Kebakaran:} Setiap perubahan kebakaran disimpan ke penyimpanan terpartisi melalui jalur \textit{transactional outbox} dengan tingkat keparahan sesuai ambang spesifikasi. Karena kontrak \texttt{HazardEvent} tidak memiliki kolom siklus hidup khusus, kebakaran yang dipensiunkan tetap tersimpan dengan severity \texttt{NORMAL} serta metadata eksplisit \texttt{attributes.lifecycle\_status = "retired"}, \texttt{retired\_reason}, \texttt{merged\_into}, dan \texttt{severity\_before\_retire}.
    \item \textbf{\textit{Checkpointing} Berkeadaan dan Koordinasi HA:} Keadaan mesin disimpan sebagai \textit{snapshot} JSON bersama \textit{offset} terakhir ke tabel \texttt{satellite\_checkpoints} pada unit koordinator \textit{shard} setiap 5~detik dan saat proses berhenti. \textit{Hazard} selalu dipersistensikan sebelum \textit{checkpoint} diperbarui, dan \textit{checkpoint} tidak pernah mundur. Hanya satu replika Aggregator pemegang \textit{lease} \texttt{source:SATELIT} yang membuka aliran SSE dan seluruh penulisannya dipagari \textit{fencing token}, sekaligus mematuhi batasan satu koneksi aktif per token kelompok pada feed satelit.
    \item \textbf{Keadaan Memori Terbatas (\textit{Bounded State}) dan Observabilitas:} Sebelum setiap \textit{checkpoint}, kebakaran yang deteksi terbarunya tertinggal melebihi \path{SATELLITE_STATE_RETENTION_HOURS} (bawaan 48~jam) dari waktu kejadian terbaru ditutup dan deteksinya dikeluarkan dari memori (\texttt{satellite\_state\_evicted}). Karena deteksi baru hanya diterima dalam jendela 6~jam \textit{watermark} dan terhubung dalam 360~menit, kebakaran yang tertinggal $>12$~jam tidak mungkin tumbuh lagi. Aggregator mengekspor metrik \textit{lag} event, \textit{offset} terapan/terbaru, histogram latensi persistensi, jumlah kebakaran aktif/singleton, deteksi terlambat, serta ukuran \textit{checkpoint} ke \texttt{/metrics} dan divisualisasikan pada dasbor Grafana (\textit{Satellite wildfire stream}).
\end{enumerate}

\vspace{1em}

\subsection{Alasan Keputusan dan Trade-Off}
Pemilihan mekanisme pengelompokan dan penyanggaan aliran satelit dievaluasi terhadap beberapa alternatif sebagaimana diringkas pada Tabel~\ref{tab:p4_tradeoff}.

\begin{footnotesize}
\begin{xltabular}{\textwidth}{p{3.5cm} c X}
\caption{Perbandingan Alternatif Pemrosesan Aliran Kebakaran Satelit (Problem 4)} \label{tab:p4_tradeoff} \\
\toprule
\textbf{Alternatif Pendekatan} & \textbf{Keputusan} & \textbf{Evaluasi dan \textit{Trade-Off}} \\
\midrule
\endfirsthead

\multicolumn{3}{l}{\textit{\footnotesize Lanjutan Tabel \ref{tab:p4_tradeoff}}} \\
\toprule
\textbf{Alternatif Pendekatan} & \textbf{Keputusan} & \textbf{Evaluasi dan \textit{Trade-Off}} \\
\midrule
\endhead

\bottomrule
\endfoot

DBSCAN per jendela tetap (\textit{tumbling window}) & Ditolak & Sederhana, tetapi memotong kebakaran yang berlangsung melintasi batas jendela dan tidak mendukung operasi \texttt{update}/\texttt{retract} secara inkremental. \\
\midrule
Pemrosesan ulang penuh pada setiap event masuk & Ditolak & Menjamin kebenaran absolut, namun kompleksitas waktu tumbuh linier terhadap seluruh riwayat aliran sehingga gagal mengejar percepatan 60$\times$. \\
\midrule
Menyangga feed satelit ke Kafka sebelum diproses & Ditolak & Redundan karena feed SSE sudah menyediakan pemutaran ulang berbasis \textit{offset} monoton; \textit{checkpoint} lokal pada penyimpanan Aggregator sudah memadai. \\
\midrule
\textbf{Komponen terhubung inkremental + indeks pita lintang + \textit{watermark} \& retensi} & \textbf{Dipilih} & Setiap event hanya memutakhirkan komponen terdampak dengan hasil identik terhadap eksekusi tanpa gangguan. \textit{Trade-off}-nya adalah perlunya \textit{snapshot checkpoint} dan kebijakan retensi 48~jam agar memori tetap terbatas. \\
\end{xltabular}
\end{footnotesize}

\vspace{1em}

\subsection{Batasan yang Disadari}
Implementasi Problem 4 memiliki beberapa batasan teknis yang disadari:
\begin{enumerate}
    \item \textbf{Pemetaan Nama Wilayah Provinsi:} Berkas referensi titik pusat provinsi BNPB belum tersedia secara bawaan di repositori; tanpa konfigurasi \path{SATELLITE_PROVINCE_CENTROIDS_FILE}, kolom \texttt{area\_name} diisi dengan label koordinat geografis.
    \item \textbf{Revisi Setelah Jendela Retensi:} Event \texttt{update} atau \texttt{retract} yang tiba setelah kebakarannya ditutup oleh batas retensi 48~jam tetap dicatat pada metrik, namun tidak lagi mengubah kebakaran yang telah ditutup.
    \item \textbf{Pengiriman Event Hilir Saat Pemulihan:} Pemutaran ulang dari \textit{checkpoint} setelah \textit{crash} dapat mengirim ulang event ke topik Kafka (\textit{at-least-once}), yang aman karena konsumen hilir melakukan deduplikasi.
    \item \textbf{Penutupan Berbasis Waktu Kejadian:} Kebijakan retensi menutup kebakaran semata berdasarkan jeda waktu kejadian dari deteksi terbaru, sehingga kebakaran yang benar-benar padam dan kebakaran yang sekadar tidak teramati satelit diperlakukan sama.
\end{enumerate}

\vspace{1em}

\subsection{Bukti Penyelesaian}
Keberhasilan Problem 4 dibuktikan melalui sesi latihan feed langsung, uji lonjakan sintetis, uji pemulihan \textit{crash}, serta pengamatan langsung pada klaster 3-node:
\begin{enumerate}
    \item \textbf{Pemrosesan Feed Live pada Klaster 3-Node (20 November 2026, \path{docs/evidence/m2/p1-p4-baseline/satellite-coastal.log}):} Dari 3 replika Aggregator aktif, hanya pemegang \textit{lease} (\texttt{aggregator-547cc65b86-nf9ld}) yang mengalirkan feed SSE hingga \textit{offset} 2.127.952 dengan \texttt{lag\_events: 0}, \texttt{p95\_persist\_ms: 0,17--0,18~ms}, 394 kebakaran aktif dalam memori, 33.560 total rekaman \texttt{WILDFIRE} terpersistensi, 0 \texttt{rejected\_events}, dan 0 \texttt{late\_beyond\_watermark}. Ketika koneksi mengalami \texttt{unexpected EOF} pada \textit{offset} 2.127.952, klien menyambung ulang dalam 500~ms dan kembali ke \texttt{lag\_events: 0}. Eviksi keadaan terbatas (\texttt{satellite\_state\_evicted}, retensi 48~jam) teramati aktif membersihkan kebakaran usang (1 kebakaran ditutup / 39--64 deteksi dikeluarkan per siklus). Dua replika cadangan tetap siaga tanpa membuka koneksi ganda.
    \item \textbf{Sesi Latihan Feed Live Lokal (17--18 November 2026):} Dengan 1 Aggregator dan 3~\textit{shard} $\times$ 2 unit PostgreSQL, \textit{backlog} 119.142 event terkejar dalam $\approx 4,5$~menit hingga \textit{lag} 0 (0 ditolak, 0 terlambat, $p95$ 3,7--5,1~ms), menghasilkan 1.872 kebakaran aktif dan 12 kebakaran pensiun yang identik antara primer dan replika tanpa satu pun deteksi ganda di dua kebakaran. Pemutusan koneksi pada \textit{offset} 112.671 dan \textit{restart} pada \textit{offset} 120.311 berhasil memulihkan 1.889 kebakaran dan kembali ke \textit{lag} 0 dalam hitungan detik.
    \item \textbf{Ketahanan Lonjakan (\textit{Burst Catch-Up}, Kriteria 4):} Pengujian feed sintetis 57.821 event yang dikirim sekaligus selesai diproses dalam 23,8~detik ($\approx 2.400$~event/detik) dan mencapai kondisi mutakhir 3,9~detik setelah pengiriman berakhir. Optimasi pelacakan progres pada PR~\#46 menghilangkan \textit{bottleneck} pembangunan ulang seluruh kebakaran per event (yang sebelumnya menurunkan laju dari 370 menjadi 40~event/detik).
    \item \textbf{Ketahanan Terhadap \textit{Crash} dan Koordinasi HA (Kriteria 5 \& 8):} Uji \texttt{TestSatelliteWorkerCrashRecoveryMatchesUninterruptedRun} menyuntikkan kegagalan pada setiap \textit{offset} penulisan dengan dua interval \textit{checkpoint}; seluruh pemulihan menghasilkan himpunan \textit{hazard} yang identik dengan eksekusi tanpa gangguan dan \textit{checkpoint} tidak pernah melampaui \textit{offset} gagal. Uji dua \textit{worker} membuktikan replika cadangan tidak mengalirkan feed saat \textit{lease} dipegang dan langsung melanjutkan dari \textit{offset} 8 setelah pemegang awal berhenti.
\end{enumerate}

\vspace{1em}

\subsection{Jawaban atas Butir ``Di Laporan''}
Berikut jawaban ringkas atas seluruh butir ``Di laporan'' spesifikasi Problem 4:
\begin{enumerate}
    \item \textbf{Identifikasi Persoalan dan Konsep Literatur:} Kriteria 2 merupakan persoalan \textit{event time} vs. \textit{processing time}, \textit{watermark}, dan \textit{allowed lateness} \cite{akidau2015dataflow,kleppmann2017designing}. Kriteria 4 merupakan \textit{backpressure} dan pengejaran \textit{backlog} (\textit{catch-up}). Kriteria 5 merupakan \textit{stateful stream processing} dengan \textit{checkpointing} yang konsisten terhadap \textit{offset} sumber.
    \item \textbf{Evaluasi Pola Pengambilan M1:} Pola \textit{polling} M1 ditinggalkan untuk sumber satelit. Feed satelit bersifat dorong (SSE) dengan \textit{offset} monoton tanpa celah; \textit{polling} berkala hanya menambah latensi dan risiko celah, sedangkan aliran SSE berbasis \textit{offset} memberikan titik lanjut yang presisi dan hemat sumber daya.
    \item \textbf{Kriteria Finalisasi Kebakaran, Penanganan Keterlambatan, dan Lokasi Penyimpanan Keadaan:} Kebakaran diterbitkan begitu memenuhi syarat ($\ge 2$ deteksi terhubung) dan diperbarui setiap ada perubahan. Kebakaran dianggap final dan ditutup dari memori setelah deteksi terbarunya tertinggal melampaui jendela retensi 48~jam dari waktu kejadian terbaru (karena tidak mungkin lagi terhubung dalam batas 360~menit maupun 6~jam \textit{watermark}). Deteksi terlambat ditangani oleh \textit{watermark} 360~menit (diterapkan bila masih dalam batas, dicatat terpisah bila melampaui). Seluruh keadaan sementara (deteksi, komponen terhubung, peta ID, \textit{watermark}, dan statistik) disimpan sebagai \textit{snapshot} JSON bersama \textit{offset}-nya di tabel \texttt{satellite\_checkpoints} pada unit koordinator \textit{shard}.
    \item \textbf{Jaminan Semantik Pemrosesan:} Sistem memberikan jaminan \textit{effectively-once} pada penyimpanan kanonis: \textit{hazard} selalu dipersistensikan sebelum \textit{checkpoint}, pemulihan pasca-\textit{crash} memutar ulang aliran tepat dari \textit{offset} \textit{checkpoint}, dan operasi \textit{upsert} bersifat idempoten (dibuktikan melalui uji \textit{crash} di setiap \textit{offset} dan pemulihan feed live). Pengiriman ke topik Kafka bersifat \textit{at-least-once} dengan deduplikasi di sisi konsumen.
\end{enumerate}
